\documentclass[10pt,a4paper]{amsart}
\usepackage[margin=19mm]{geometry}
\usepackage{amsmath,amssymb,mathtools}
\usepackage[scr=rsfs,cal=euler]{mathalfa}
\usepackage{xfrac}
\usepackage[unicode,hidelinks,bookmarks=false]{hyperref}
\newcommand{\ie}{i.e.\ }
\newcommand{\cf}{cf.\ }
\newcommand{\resp}{respectively\ }
\newcommand{\Subsection}{\subsection}
\providecommand{\nobreakdash}{\nobreak\hbox{-}}
\setlength{\emergencystretch}{2em}
\multlinegap0pt
\usepackage{bm}
\usepackage{bbm}
\usepackage{dsfont}
\usepackage{xspace}
\usepackage{cancel}
\usepackage{mathtools}
\usepackage{amsthm}
\makeatletter
\@ifundefined{theorem}{\newtheorem{theorem}{Theorem}}{}
\@ifundefined{lemma}{\newtheorem{lemma}[theorem]{Lemma}}{}
\makeatother
\title[Noetherianity of polynomial rings up to group actions]{Noetherianity of polynomial rings up to group actions}
\author{Liping Li, Yinhe Peng,, Zhengjun Yuan}
\date{Source: arXiv:2502.14306v2 (CC BY 4.0)}
\begin{document}
\maketitle

\noindent\emph{Source and licence.} Noetherianity of polynomial rings up to group actions. Version arXiv:2502.14306v2 (CC BY 4.0), distributed under the Creative Commons Attribution 4.0 International licence, as recorded in the arXiv licence metadata for this exact version and shown on the abstract page https://arxiv.org/abs/2502.14306v2. Full rights notice: licenses/arxiv-2502.14306-CC-BY-4.0.txt.\par\smallskip

\noindent\textit{Editorial scope.} Counted unit: the Lemma whose source label is \texttt{unbounded and dense} in the article (source0.tex lines 603--612, source label \texttt{unbounded and dense}), delivered together with the article's own complete original proof of it (source0.tex lines 614--628, one \texttt{proof} environment of 15 lines). No mathematics is written, repaired or re-derived: statement and proof are the article's own text line for line. Required context retained uncounted: none. Every definition, display and label of those lines is retained unchanged. Omitted from this package and not redistributed: 1--602, 613--613, 629--870. Nothing inside a retained excerpt is omitted or altered: every retained body is delivered exactly as the article's lines are written, so no omission receipt of any kind accompanies this package. Citations: the retained excerpts carry no citation command at all, so no bibliography entry of the article is transcribed and no reference is redistributed. Reference adaptation: none was needed and none was made; every reference inside the delivered unit resolves inside it, so no unresolved reference or citation is produced. Printed slips kept literal: no printed word, symbol, hypothesis or number of the counted statement or of its proof was altered. Layout and class: the article's own class is \texttt{amsart} with packages including amssymb, amsthm, amsmath, amsfonts, graphics, hyperref, bbm, xy, enumerate, eucal, thmtools, xcolor, ytableau; the wrapper is the lane's pinned \texttt{amsart} editorial wrapper at 10pt on A4, which restates the theorem-like environments the retained text needs and adds the article's own macro definitions that the retained text needs (none), with extra wrapper packages (bm, bbm, dsfont, xspace, cancel, mathtools). The delivered numbering is the unit's own. Assets: no retained span contains a figure environment, a graphics-inclusion command, a PDF-page inclusion, a table or a TikZ picture; no image file is copied into this package, the package declares no asset dependency and no image file is redistributed with it. Licence: version arXiv:2502.14306v2 of this article is distributed under the Creative Commons Attribution 4.0 International licence, as recorded in the arXiv licence metadata for this exact version and shown on the abstract page https://arxiv.org/abs/2502.14306v2. A full rights notice is delivered at licenses/arxiv-2502.14306-CC-BY-4.0.txt. Characters: every retained excerpt is pure ASCII, so the delivered engine, XeLaTeX with its default Latin Modern text font, reports no missing character.
\begin{lemma} \label{unbounded and dense}
Let $\leqslant$ be a linear order on $S$.
\begin{enumerate}
\item If $\leqslant$ is homogenous, then it is dense and unbounded.

\item If $\leqslant$ is unbounded, then it is homogenous if and only if for $a < b$ and $a' < b$ in $S$, one has $(a, b) \cong (a', b')$.

\item The linear order $\leqslant$ is globally homogenous if and only if for every $a < b$ in $S$, one has $(a, b) \cong (S, \leqslant)$.
\end{enumerate}
\end{lemma}

\begin{proof}
(1). If $(S, \leqslant)$ has an endpoint, say, a minimal element $a$. Then for any $b \neq a$ in $S$, the open interval $(-, a)$ is empty, while $(-, b)$ is nonempty, so they cannot be isomorphic. Thus $(S, \leqslant)$ is unbounded. Similarly, if $S$ is not dense, then we can find an increasing pair $a < b$ in $S$ with $(a, b) = \emptyset$. Since $S$ is unbounded, we can find another $c \in S$ with $c > b$. Then $(a, b)$ is not isomorphic to $(a, c)$. This contradiction tells us that $(S, \leqslant)$ is dense.

(2). The only if direction is clear. For the other direction, since $\leqslant$ is unbounded, we can find a certain $c \in S$ such that $c > a$ and $c > a'$. Then one has $(a, c) \cong (a', c)$ and hence
\[
(a, -) = (a, c) \sqcup [c, -) \cong (a', c) \sqcup [c, -) = (a', -).
\]
Similarly, one can construct a desired isomorphism $(-, a) \cong (-, a')$.

(3). The only if direction is clear. For the if direction, one needs to show $(-, a) \cong (S, \leqslant) \cong (a, -)$ for $a \in S$. We prove the first isomorphism since the second one can be verified similarly.

We claim that $\leqslant$ is unbounded. Otherwise, suppose that $s_0$ is the minimal element in $S$. For any $s < t$ in $S$, let $\sigma: (s, t) \to (S, \leqslant)$ be an isomorphism. Then one can find a certain $s' \in (s, t)$ such that $\sigma(s') = s_0$. Consequently, one has $s' > s$ and $s'$ is the minimal element in $(s, t)$. But then $(s, s')$ is the empty set, which cannot be isomorphic to $(S, \leqslant)$. This contradiction tells us that $S$ has no minimal element. Similarly, one can show that it has no maximal element.

Choose an element $x \in (s, t)$. Then the isomorphism $\sigma: (s, t) \to (S, \leqslant)$ induces an isomorphism $(s, x) \cong (-, \sigma(x))$. But $(s, x) \cong (S, \leqslant)$ as well, so it remains to show that $(-, a) \cong (-, \sigma(x))$. But this is clear since by statement (2), $\leqslant$ is homogenous.
\end{proof}

\end{document}
